\documentclass[10pt]{article}
\usepackage{parskip}
\usepackage[margin=1in]{geometry} 
\usepackage{tikz}
\usetikzlibrary{positioning}
\usepackage{amsmath,amsthm,amssymb, graphicx, multicol, array}
\usepackage{enumitem}
\usepackage{amssymb}
\usepackage{float}
\usepackage{href-ul}
\usepackage[bb=boondox,bbscaled=.95,cal=boondoxo]{mathalfa} 

 
\newcommand{\N}{\mathbb{N}}
\newcommand{\Z}{\mathbb{Z}}
 
\newenvironment{problem}[2][Problem]{\begin{trivlist}
\item[\hskip \labelsep {\bfseries #1}\hskip \labelsep {\bfseries #2.}]}{\end{trivlist}}

\begin{document}
 
\title{Experimental design --- Portfolio task 2}
\author{Jakob Sverre Alexandersen\\
GRA4158 Marketing and the Analysis of Experiments and Quasi-Experiments}
\maketitle

Your task is to propose a marketing-relevant research question (business problem) and discuss
different design choices to best test the treatment effect in an experiment. Clearly specify the
independent and dependent variables and identify potential challenges in measurement and
manipulation of the focal variables (e.g., how to operationalize the treatment and
counterfactual). Discuss pros and cons of different experimental settings (lab or field; between
or within) and threats to internal and external validity of these design choices for your questions
and select the best fitting design for your business problem / research question.

\textbf{Assessment dimensions}
\begin{enumerate}
    \item Research question
    \begin{itemize}
        \item Identifies a relevant and non-trivial research question / business problem (the question
        does not need to be new in the sense of identifying an academically new research idea,
        but it should be practically relevant and not trivial to answer).
        \item Clearly describes the problem and focal treatment effect(s) to be tested.
    \end{itemize}
    \item Operationalization and variables
    \begin{itemize}
        \item Clearly specifies the theoretical independent variable (what is varied) and dependent
        variable (focal outcome) and appropriately distinguishes between these and actual
        manipulation and measurement of the outcome.
        \item Provides a concrete description of the treatment condition and the counterfactual
        (control group) and why this counterfactual is most appropriate for the research
        question.
        \item Reflects on potential challenges in manipulating the IV and measuring the DV.
    \end{itemize}
    \item Experimental design strategy
    \begin{itemize}
        \item Provides justified arguments for the settings (lab vs. field), explicitly discussing the trade-
        offs between control (internal validity) and realism (external validity) relative to the
        selected research question / business problem.
        \item Explicitly chooses a design structure (between-subjects vs. within-subjects) and defends
        this choice based on the nature of research question and concerns about validity threats
    \end{itemize}
    \item Validity Analysis
    \begin{itemize}
        \item Comprehensively discusses specific threats to internal validity and external validity
        applicable to the considered designs choices and research question, as discussed in the
        course.
        \item Goes beyond simple explanations of threats to discuss how the proposed design
        mitigates these specific threats (or acknowledges which threats remain unavoidable).
    \end{itemize}
    \item Use of Terminology and Concepts:
    \begin{itemize}
        \item Accurately applies concepts and terminology from the class or course literature
    \end{itemize}
    \item Clarity and Coherence
    \begin{itemize}
        \item Provides a logical structure and coherent organization of the document.
        \item Communicates findings and insights clearly and effectively, using concise and well-structured language
    \end{itemize}
    \item Academic Rigor
    \begin{itemize}
        \item Demonstrates critical thinking and deep analysis, explaining the "why" behind decisions
        rather than just describing them.
        \item Accurately applies specific course concepts and terminology.
        \item References relevant literature (course readings or to a limited degree additional high-quality academic literature where necessary) to support arguments.
    \end{itemize}
\end{enumerate}
\textbf{Format}
\begin{enumerate}
    \item Written report (not bullet points)
\end{enumerate}
\textbf{Submission}
\begin{itemize}
    \item The submission should not have much more than 2500 words (excluding references,
    tables, and figures).
    \item During the course you might get additional ideas to supplement your design. You can
    always modify and extend your discussion until the final deadline.
    \item Integrate your final report into the portfolio as one well-organized document.
    \item Submit the final version as part of the portfolio in Wiseflow.
\end{itemize}

\newpage 

\textbf{Does Authenticity Sell? Designing an Experiment to Test the Effect of Influencer Content Style on Purchase Intent} 

\section{Introduction and Research Question}

The relationship between younger consumers and advertising is undergoing a fundamental shift. Adoption of ad-blocking software among young internet users has grown at approximately 40\% year-over-year, and eye-tracking research on digital natives reveals that intrusive advertising formats (pop-ups, non-skippable video ads, and content-reorganizing banners) are perceived as deeply annoying and are actively avoided [CITE{futurenet2024}]. This is not merely a technological trend; it reflects a generational change in how advertising credibility is evaluated. [CITE{djafarova2017}] found that young consumers perceive non-traditional content creators, like bloggers, and YouTubers, as significantly more credible and persuasive than traditional celebrity endorsers or conventional brand advertisements. More recent work confirms that influencer authenticity has a strong positive effect on Gen~Z trust and purchase intentions [CITE{acr2025}], and that Gen~Z consumers view authentic influencers as ``highly educated friends'' from whom they seek advice, rather than as paid advertisers [CITE{djafarova2024}].

Given this shift, brands are investing heavily in influencer marketing. However, the question of \textit{what style} of influencer content is most effective remains open. On one end of the spectrum, brands commission polished, high-production influencer partnerships with scripted copy and studio-quality visuals. On the other end, a growing number of creators adopt a raw, casual aesthetic: filming on their phone, speaking without a script, and using humor that borders on self-deprecating (e.g., ``honestly just buy this, I don't know what else to tell you''). Observational data might suggest that authentic-looking content generates higher engagement, but such observations are heavily confounded: platforms' targeting algorithms amplify content to audiences predicted to engage with it, and users self-select into following creators whose style they prefer. As [CITE{braun2020}] demonstrate, even randomized A/B tests conducted through platform-provided tools suffer from \textit{divergent delivery}, where the algorithm serves different ad variants to systematically different user populations, making it impossible to attribute differences to content alone.

This report proposes an experiment to test the following causal question: \textbf{does casual, authentic-style influencer content cause higher purchase intent than professionally produced influencer content?} The focal treatment effect concerns the \textit{perceived authenticity} of the content style, not the presence or absence of influencer marketing itself. The remainder of this report specifies the independent and dependent variables, discusses the operationalization of the treatment and counterfactual, evaluates alternative experimental designs, and analyzes threats to internal and external validity.

\section{Variables and Operationalization}

\subsection{Independent Variable}

The theoretical independent variable is content authenticity style; the degree to which influencer content appears spontaneous, unpolished, and genuine versus professionally produced and scripted. This construct is operationalized through two experimental conditions:

The treatment condition presents casual, authentic influencer content. Concretely, this is a short video (approximately 30--60 seconds) filmed on a smartphone in a natural setting (e.g., the influencer's home), with informal language, minimal editing, and a conversational tone. The influencer speaks directly to the camera without a visible script, and the product endorsement is embedded in casual commentary.

The control condition presents professional, polished influencer content. This is a video of identical length featuring \textit{the same influencer} promoting \textit{the same product}, but filmed with studio lighting, high-quality audio, scripted copy, and polished editing resembling a traditional sponsored post.

\subsection{Counterfactual and the Compound Treatment Problem}

The choice to use the same influencer and the same product across both conditions is a deliberate design decision to avoid the \textit{compound treatment problem}. As discussed in the course, a compound treatment arises when the experimental manipulation unintentionally varies multiple factors simultaneously. If we compared influencer~A (casual) with influencer~B (professional), any observed difference could be attributed to the person, their audience, their delivery style, or the content format, and not solely to authenticity. By holding the influencer and product constant, the manipulation isolates the content style dimension as cleanly as possible.

The counterfactual is an \textit{explicit alternative} rather than the absence of content. This is appropriate because the research question concerns the \textit{relative} effectiveness of two content strategies, not whether influencer content works at all. A ``no content'' control group would answer a different question and would not provide actionable guidance for the business decision at hand.

\subsection{Dependent Variable}

The primary dependent variable is purchase intent, which captures the consumer's likelihood of purchasing the promoted product after viewing the content. In the proposed field setting, this is operationalized through behavioral measures recorded via a unique tracked link embedded in each post. Each experimental condition (casual and professional) contains a distinct UTM-tagged or affiliate URL pointing to the product page, allowing the researcher to attribute downstream behavior---click-through, add-to-cart, and conversion---to the specific content variant that generated it. These revealed-preference measures avoid the well-documented gap between stated intentions and actual purchasing behavior, providing a direct read on how content style influences real consumer decisions. Were the experiment instead conducted in a lab, the dependent variable would be measured using a validated 7-point Likert scale (e.g., ``How likely are you to purchase this product?'') supplemented by a hypothetical willingness-to-pay measure---useful proxies, but inherently less actionable than observed behavior.

\subsection{Challenges in Manipulation and Measurement}

Two key challenges deserve attention. First, regarding manipulation effectiveness: casual content may vary not only perceived authenticity but also perceived humor, effort, or production quality. In a lab setting, this would be verified through explicit manipulation checks where participants rate perceived authenticity, humor, and production quality. In the proposed field setting, however, such self-report measures are unavailable. The researcher cannot directly confirm that authenticity is the primary dimension being varied, which means the field experiment tests the \textit{bundle} of features that distinguish casual from professional content rather than authenticity in isolation. This is a genuine limitation of the field design and one of the strongest arguments for a complementary lab study to probe the underlying mechanism.

Second, regarding measurement of the DV: although the tracked link provides a clean attribution path from content to behavior, the resulting measures are noisy. A single click does not imply genuine purchase intent, and conversion rates in influencer campaigns are typically low, requiring large sample sizes to detect modest effects. The link also introduces a minor demand artifact: users who notice the tracked URL may behave differently than they would with an organic link, though this concern is attenuated by the fact that UTM parameters and affiliate links are standard practice in influencer content and are unlikely to be salient to most users. Furthermore, behavioral data alone cannot reveal \textit{why} participants acted as they did---the mechanism remains a black box. A lab experiment would allow the researcher to supplement outcome measures with mediator scales (e.g., perceived authenticity, trust), but this diagnostic depth comes at the cost of ecological validity. The field design accepts this trade-off, relying on the theoretical framework to support mechanistic interpretation.

\section{Experimental Design Strategy}

\subsection{Between-Subjects Design}

The proposed experiment uses a between-subjects design, in which each participant is randomly assigned to view either the casual or the professional video, but not both. This choice is driven by the nature of the treatment. A within-subjects design, where the same participant views both videos, would be inappropriate for two reasons. First, it would introduce severe \textit{carryover effects}: once a participant has seen the casual version, their reaction to the professional version (or vice versa) would be contaminated by the comparison. Second, it would create strong \textit{demand effects}, as participants would almost certainly infer that the study is comparing the two styles and adjust their responses accordingly. Both threats are discussed extensively in this course as primary concerns for within-subjects designs.

Between-subjects designs eliminate these concerns at the cost of requiring more participants (since each participant provides data for only one condition) and introducing the possibility of pre-existing differences between groups. However, the latter concern is addressed through random assignment, which ensures that (in theory) the two groups are equivalent on all observed and unobserved characteristics. To further strengthen balance on key demographics (age, gender, social media usage frequency), \textit{stratified random assignment} can be employed, randomizing within homogeneous subgroups rather than across the full sample. This is particularly important if we wish to analyze heterogeneous treatment effects across demographic subgroups.

\subsection{Lab vs.\ Field Experiment}

\subsubsection{The Case for a Lab Experiment}

An online lab experiment would offer maximum control over the experimental setting. Each participant could be randomly assigned to view exactly one video within a survey interface, eliminating any interference from platform algorithms. The researcher could include manipulation checks (e.g., ``How authentic did this content feel?'') and mediator scales (perceived trust, perceived effort) alongside the outcome measures. This diagnostic depth is the core advantage of a lab setting: it allows the researcher not only to estimate the treatment effect but also to probe the mechanism through which it operates. The Hawthorne effect is reduced in an online lab relative to a physical lab, since participants complete the task from their own environment.

However, the lab setting carries substantial limitations for this particular research question. Watching an influencer video inside a survey is fundamentally different from encountering it in a social media feed. The social context (likes, comments, share counts, algorithmic placement among other content) is entirely absent, and these cues may moderate the authenticity effect in practice. Casual content may derive part of its appeal from the \textit{contrast} with polished content in the same feed, an effect that disappears in isolation. Furthermore, the dependent variable in a lab must rely on stated purchase intent rather than actual behavior, introducing the well-documented intention--behavior gap. For a research question whose ultimate business relevance lies in real purchasing outcomes, this is a significant sacrifice.

\subsubsection{The Case for a Field Experiment}

A field experiment, for instance, running an A/B test through Instagram or TikTok's advertising tools, would score highly on the four key features of field experiments identified in this course: real-world context, authenticity of treatment, representativeness of participants, and relevant outcome measures. Participants would encounter the content organically in their social media feed, respond with real behavior (clicks, purchases), and the sample would reflect the actual target population. The ecological validity of such a design would be high, and the dependent variable could be operationalized as actual click-through, add-to-cart, or conversion. These are behavioral measures that directly answer the business question.

The primary internal validity threat of a field experiment is \textit{divergent delivery}. As \cite{braun2020} demonstrate, platform A/B testing tools optimize ad delivery based on predicted user-content relevance, meaning that the casual video and the professional video may be served to systematically different user populations. \cite{gordon2019} reinforce this concern by showing that even with rich individual-level data, selection effects in digital advertising cannot be adequately controlled for post hoc. Additionally, a field experiment does not easily permit manipulation checks, so the researcher cannot directly verify that the treatment operates through perceived authenticity rather than some other channel. Non-compliance is also a concern: not all users assigned to see an ad will actually be exposed, creating one-sided non-compliance that shifts the estimand from ATE to ITT or LATE.

\subsubsection{Recommended Design: Field Experiment with Mitigation Strategies}

Despite the internal validity threats, the proposed design is a field experiment conducted through a social media platform's advertising infrastructure. The core justification is that the research question is inherently about how consumers respond to influencer content \textit{in its natural environment}. Content style effects are likely moderated by the social media context itself; the feed, the social proof signals, the fleeting attention span of scrolling users, and stripping away that context in a lab would risk studying a phenomenon that does not transfer to the real world. As \cite{gordon2019} note, the value of field experiments in digital advertising lies precisely in their ability to measure effects on real behavior in real markets.

To address the divergent delivery problem, several mitigation strategies can be employed. First, the brand can configure the platform's ad manager to use \textit{identical audience targeting} for both conditions, restricting delivery to a pre-defined demographic segment (e.g., Instagram users aged 18--30 in a specific geographic region) and disabling algorithmic optimization of creative delivery where the platform allows it. Second, each content variant includes a unique tracked link (UTM-tagged or affiliate URL) to the product page, enabling clean attribution of click-through, add-to-cart, and conversion behavior to the specific condition that generated it. Third, the analysis should report both the intent-to-treat (ITT) estimate and, where exposure data is available, the local average treatment effect (LATE) to account for non-compliance. Fourth, demographic and behavioral covariates collected by the platform (age, gender, usage frequency) can be included in the regression to improve precision and check for balance across conditions. While these measures cannot fully eliminate the divergent delivery concern, they substantially reduce it and make the remaining threat transparent and quantifiable.

The sacrifice relative to a lab experiment is the inability to include manipulation checks or mediator scales. The researcher cannot directly verify that any observed difference operates through perceived authenticity. This limitation should motivate a complementary lab study as a follow-up to probe the mechanism once the field experiment establishes whether a practically meaningful effect exists in the real market.

\section{Validity Analysis}

\subsection{Threats to Internal Validity}

The most serious internal validity threat is that the platform's algorithm may serve the two ad variants to systematically different audiences, even when targeting parameters are held constant. As \cite{braun2020} document, this \textit{divergent delivery} means that any observed difference in outcomes could reflect audience composition rather than content style. The mitigation strategies discussed above, restricting targeting, disabling creative optimization where possible, and checking covariate balance post hoc reduce but do not fully eliminate this threat. Transparency about residual confounding is essential when interpreting the results.

In a field experiment, not all users assigned to see an ad will actually view it. Some will scroll past before the video loads, others may be served the ad but not attend to it. This one-sided non-compliance means the experiment estimates an intent-to-treat (ITT) effect rather than the average treatment effect on the treated (ATT). Where the platform provides impression-level data, an instrumental-variables approach can recover the local average treatment effect (LATE), but the ITT remains the most policy-relevant estimand for managerial decision-making, since it reflects what happens when the brand \textit{deploys} a given content strategy.

Despite using the same influencer and product across conditions, the casual and professional videos inevitably differ on multiple dimensions, not only authenticity but also humor, perceived effort, visual aesthetics, and potentially video pacing. This means the treatment is not a clean manipulation of a single construct. Unlike in a lab setting, the field experiment does not allow manipulation checks to diagnose which dimensions are driving the effect. This limitation is inherent to the field design and should be acknowledged: the experiment tests the \textit{bundle} of features that constitute casual versus professional content, not authenticity in isolation. Including platform-provided demographic and behavioral covariates in the analysis can partially account for heterogeneous responses, but it cannot decompose the treatment into its constituent mechanisms.

If the same user follows the influencer organically and encounters both the casual and professional content (e.g., through the influencer's own feed alongside the paid ad), treatment contamination may occur. This risk can be reduced by running the two ad variants in non-overlapping time windows or geographic regions, though this introduces potential confounds from temporal or regional differences. Alternatively, the brand can coordinate with the influencer to ensure that only the paid ads are posted during the experimental window.

Behavioral outcomes such as conversion are binary and typically have low base rates in influencer campaigns, meaning the experiment requires a large sample to detect modest effects. A formal power analysis based on historical conversion rates for similar campaigns should be conducted prior to launch. Under-powering the experiment risks a Type~II error that incorrectly concludes content style does not matter.

\subsection{Threats to External Validity}

The norms and expectations of content differ across platforms. Casual, ``off-the-cuff'' content is native to TikTok but may feel out of place on LinkedIn or even Instagram. The field experiment is conducted on a single platform, and results should not be generalized to other platforms without replication. However, running the experiment \textit{on} a platform, rather than abstracting away the platform context in a lab, ensures that the estimated effect is valid for at least one real market environment.

The effect of content style likely depends on the product category. Casual, humorous content may be highly effective for ``enjoyment-driven'', low-involvement products (e.g., snacks, fashion accessories) but less so for utilitarian, high-involvement products (e.g., financial services, electronics). The proposed experiment tests a single product category, and the results should not be generalized beyond it without replication.

The effect may depend on the particular influencer's persona. An influencer whose brand is already casual may show a smaller treatment effect than one whose audience expects polished content. While using a single influencer controls for person-level confounds, it limits generalizability. Future replications with multiple influencers would strengthen external validity.


Social media trends evolve rapidly. What reads as ``authentic'' today may feel forced or faux\footnote{Like Wendy's replies on $\mathbb{X}$ (formerly Twitter)} in six months as the casual aesthetic becomes mainstream. The field experiment captures a snapshot of consumer preferences at a particular moment, and its conclusions may have a limited shelf life.

\section{Discussion and Conclusion}

This report has proposed an experimental design to test whether casual, authentic-style influencer content causes higher purchase intent than professionally produced influencer content. The focal treatment effect concerns perceived authenticity of the content style, operationalized through two video conditions featuring the same influencer and the same product but differing in production quality, tone, and delivery.

The recommended design is a between-subjects field experiment conducted through a social media platform's advertising infrastructure. This choice prioritizes external validity and behavioral realism; a trade-off that is justified by the nature of the research question. Content style effects are deeply embedded in the social media context: the feed environment, the social proof cues, and the fleeting attention of scrolling users all shape how authenticity is perceived and whether it translates into action. A lab experiment, while offering superior control and the ability to probe mechanisms through manipulation checks, would strip away precisely the contextual factors that make the question business-relevant in the first place.

The field design is not without limitations. The divergent delivery problem documented by \cite{braun2020} means that platform algorithms may confound the treatment comparison by serving the two ad variants to systematically different audiences. \cite{gordon2019} further demonstrate that observational corrections for such selection effects are unreliable. The mitigation strategies proposed like restricting targeting parameters, disabling creative optimization, and checking covariate balance, reduce but do not eliminate this threat. Furthermore, the compound treatment concern remains: without manipulation checks, the experiment cannot isolate perceived authenticity from the bundle of features that distinguish casual from professional content. These limitations should motivate a complementary lab study as a follow-up to probe the mechanism once the field experiment establishes whether a practically meaningful effect exists in the real market.

From a practical standpoint, the answer to this research question has direct implications for how brands allocate their influencer marketing budgets. If casual content proves more effective at driving real conversions, the irony is notable: the less brands invest in production, the more they may earn in return \dots at least among the younger consumers who increasingly dictate market trends.

\newpage
\begin{thebibliography}{9}

\bibitem[Braun and Schwartz, 2020]{braun2020}
Braun, M. and Schwartz, E.~M. (2020).
\newblock Where A/B testing goes wrong: How divergent delivery affects what online experiments cannot (and can) tell you about how customers respond to advertising.
\newblock \textit{Working Paper, Harvard Business School}.

\bibitem[Djafarova and Rushworth, 2017]{djafarova2017}
Djafarova, E. and Rushworth, C. (2017).
\newblock Exploring the credibility of online celebrities' Instagram profiles in influencing the purchase decisions of young female users.
\newblock \textit{Computers in Human Behavior}, 68:1--7.

\bibitem[Djafarova et~al., 2024]{djafarova2024}
Djafarova, E., Kramer, T., and Schoenmueller, V. (2024).
\newblock Differences in perceived influencer authenticity: A comparison of Gen~Z and Millennials' definitions of influencer authenticity during the de-influencer movement.
\newblock \textit{Online Media and Global Communication}, 3(1).

\bibitem[Gordon et~al., 2019]{gordon2019}
Gordon, B.~R., Zettelmeyer, F., Bhatt, N., and Lovett, M.~J. (2019).
\newblock A comparison of approaches to advertising measurement: Evidence from big field experiments at Facebook.
\newblock \textit{Marketing Science}, 38(2):193--225.

\bibitem[{Advances in Consumer Research}, 2025]{acr2025}
{Advances in Consumer Research} (2025).
\newblock The impact of influencer authenticity on purchase intentions among Gen~Z consumers.
\newblock \textit{Advances in Consumer Research}, 53.

\bibitem[Pla\v{s}ilov\'{a} et~al., 2024]{futurenet2024}
Pla\v{s}ilov\'{a}, A., \v{S}krab\'{a}kov\'{a}, B., and Kita, P. (2024).
\newblock Perspectives of young digital natives on digital marketing: Exploring annoyance and effectiveness with eye-tracking analysis.
\newblock \textit{Future Internet}, 16(4):125.

\end{thebibliography}

\end{document}